\section{Extended limitations and ethics}
\label{sec:s-limitations-v2}

\subsection{Dataset and task scope}

The evidence covers two public real graphs and four measured regimes from one synthetic IBM generator. Prediction units differ across transaction nodes, account nodes, and transaction edges. No result estimates performance for an unobserved institution, jurisdiction, payment network, fraud typology, or laundering adaptation. DGraphFin time is an incident-edge-derived bucket construction, Elliptic uses fixed dataset steps, and IBM HI/LI are synthetic regimes rather than bank populations.

\subsection{Architecture and construct scope}

The primary real-data grid contains MLP, GCN, and GraphSAGE rather than a broad temporal, heterogeneous, or transformer survey. Strict-to-isolated changes graph access and may change several structural properties at once; it shows sensitivity, not a unique mechanism. V24 stress labels are construct-invalid as separate conditions, and the sender-receiver row is an identity alias. The IBM early-to-late classifier uses first-50\% labels with first-60\% label-free histories, so it is not a pure first-50\%-information simulation.

\subsection{Statistical scope}

Ten seeds quantify optimization variation on fixed datasets and splits, not uncertainty across institutions or future periods. Percentile intervals and Wilcoxon tests inherit that limit. Holm controls the declared analysis families, not every historical exploration. The IBM construction analysis averages four fixed contexts within seed; the 208 context rows remain available because dependence-aware aggregation does not erase heterogeneity. GraphSafe averages six contexts within seed and therefore gives a valid ten-block comparison at the cost of hiding some context variation in its compact summary.

\subsection{Compute and missing cells}

IBM Large, Medium GINE, fixed DGraphFin GAT h64/l2, and the max-pool extension remain unmeasured. Another implementation or larger GPU could make them feasible, but no predictive direction follows. Runtime is runner elapsed time within the recorded environment, not a general energy, latency, or hardware-efficiency benchmark.

\subsection{Artifact and validator scope}

The rebuild verifies saved-output reconstruction, not clean-machine retraining. The support schema and 14 mutations were designed inside this project and establish implementation conformance only. External curators, independent implementations, and inter-rater studies are absent. The artifact must remain synchronized with the final PDFs, but checksums cannot prove construct validity or scientific truth.

\subsection{Operational and fairness limits}

The fixed 1:5 cost, review budgets, and calibration bins are scenarios. Real investigator capacity, delayed labels, investigation quality, interventions, and feedback loops can change value~\cite{siddiqui2018feedback}. The datasets do not provide an appropriate surface for a complete fairness or disparate-impact evaluation. No claim is made about case resolution, investigator productivity, monetary loss reduction, regulatory compliance, or equitable deployment.

\subsection{Intended use and harm}

Fraud scores can burden investigators and customers through false alerts and can affect access to financial services. The benchmark evaluates ranking and review support; it does not authorize automated account closure, punitive action, or identity inference. Deployment would require institution-specific validation, governance, appeal processes, monitoring, security review, and human oversight. Dataset, model, and deployment cards should preserve intended use, exclusions, provenance, and limitations~\cite{gebru2021datasheets,mitchell2019modelcards,pushkarna2022datacards}.

\subsection{Remaining human review}

Professor/coauthor scientific review, author metadata, venue declarations, biographies if required, and the current TKDE page and supplementary-material policy remain human gates. Publication-design completion does not remove the empirical limitations above.
